% Author: Prof. Dr. Matthias Jung, DL9MJ
% Year: 2026
%
% Vergleich der Grenzwerte nach ETSI und AFuV 33/2007 (2)
% Angaben in dBc

\begin{tikzpicture}

\begin{axis}[
    axis on top,
    height=5cm,
    scale only axis,
    width=\linewidth*\getDarcImageFactor,
    xmin=0.15,
    xmax=3000,
    ymin=-75,
    log ticks with fixed point,
    xmode=log,
    xlabel={Frequenz in \unit{\mega\hertz}},
    ylabel={Grenzwert in dBc},
    %xtick={0.15,1.5,15,150,1500,3000},
    %xticklabels={
    %    $0{,}15$,
    %    $1{,}5$,
    %    $15$,
    %    $150$,
    %    $1500$,
    %    $3000$
    %},
    ytick={-75,-70,-65,-60,-55,-50,-45,-40,-35,-30},
    %grid=major,
    %major grid style={
    %    draw=DARCgray!35,
    %    line width=0.4pt
    %},
    legend style={
        at={(0.5,-0.2)},
        anchor=north,
        draw=none,
        legend columns=4,
        font=\footnotesize,
        /tikz/every even column/.append style={column sep=0.1cm}
    }
]

% AFU Bänder:

%2,2km
\fill[color=DARCgray!50!white](axis cs: 0.1357,-100) rectangle (0.1378,100);
%630m
\fill[color=DARCgray!50!white](axis cs: 0.472,-100) rectangle (0.479,100);
%160m
\fill[color=DARCgray!50!white](axis cs: 1.81,-100) rectangle (2.0,-100);
%80m
\fill[color=DARCgray!50!white](axis cs: 3.50,-100) rectangle (3.8,100);
%60m
\fill[color=DARCgray!50!white](axis cs: 5.3515,-100) rectangle (5.3665,100);
%40m
\fill[color=DARCgray!50!white](axis cs: 7.0,-100) rectangle (7.2,100);
%30m
\fill[color=DARCgray!50!white](axis cs: 10.1,-100) rectangle (10.15,100);
%20m
\fill[color=DARCgray!50!white](axis cs: 14.0,-100) rectangle (14.35,100);
%17m
\fill[color=DARCgray!50!white](axis cs: 18.068,-100) rectangle (18.168,100);
%15m
\fill[color=DARCgray!50!white](axis cs: 21.0,-100) rectangle (21.45,100);
%12m
\fill[color=DARCgray!50!white](axis cs: 24.89,-100) rectangle (24.99,100);
%10m
\fill[color=DARCgray!50!white](axis cs: 28.0,-100) rectangle (29.7,100);
%6m
\fill[color=DARCgray!50!white](axis cs: 50.0,-100) rectangle (52.0,100);
%4m
\fill[color=DARCgray!50!white](axis cs: 70.0,-100) rectangle (70.5,100);
%2m
\fill[color=DARCgray!50!white](axis cs: 144.0,-100) rectangle (146.0,100);
%70cm
\fill[color=DARCgray!50!white](axis cs: 430.0,-100) rectangle (440.0,100);
%23cm
\fill[color=DARCgray!50!white](axis cs: 1240.0,-100) rectangle (1300.0,100);
%13cm
\fill[color=DARCgray!50!white](axis cs: 2320.0,-100) rectangle (2450.0,100);

% --------------------------------------------------------------------------
% ETSI, 5 W
%
% -(43 + 10 log10(5 W)) = -49.9897 dBc
% Im Diagramm sinnvoll auf -50 dBc gerundet.
% --------------------------------------------------------------------------

\addplot[
    DARCgreen,
    very thick
]
coordinates {
    (0.15,-50)
    (3000,-50)
};
\addlegendentry{ETSI: \qty{5}{\watt}}


% --------------------------------------------------------------------------
% ETSI, 100 W
%
% Rechnerisch:
% -(43 + 10 log10(100 W)) = -63 dBc
%
% Unterhalb 30 MHz gilt laut Tabelle jedoch der absolute Grenzwert
% von -50 dBc.
% --------------------------------------------------------------------------

\addplot[
    DARCorange,
    very thick
]
coordinates {
    (0.15,-50)
    (30,-50)
    (30,-63)
    (3000,-63)
};
\addlegendentry{ETSI: \qty{100}{\watt}}


% --------------------------------------------------------------------------
% ETSI, 750 W
%
% Rechnerisch:
% -(43 + 10 log10(750 W)) = -71.75 dBc
%
% Unterhalb 30 MHz: -50 dBc
% Ab 30 MHz:        -70 dBc
% --------------------------------------------------------------------------

\addplot[
    DARCred,
    very thick
]
coordinates {
    (0.15,-50)
    (30,-50)
    (30,-70)
    (3000,-70)
};
\addlegendentry{ETSI: \qty{750}{\watt}}


% --------------------------------------------------------------------------
% Vfg. 33/2007 (2)
%
% 0.15 ... 1.7 MHz:  -60 dBc
% 1.7  ... 35 MHz:   -40 dBc
%
% 35 ... 50 MHz:
%
%   -(40 + 129.1 log10(f / 35))
%
% 50 ... 1000 MHz:   -60 dBc
% 1000 ... 3000 MHz: -50 dBc
% --------------------------------------------------------------------------

\addplot[
    DARCblue,
    very thick
]
coordinates {
    (0.15,-60)
    (1.7,-60)
    (1.7,-40)
    (35,-40)
};
\addlegendentry{Vfg. 33/2007 (2)}

% Übergang von 35 MHz bis 50 MHz
\addplot[
    DARCblue,
    very thick,
    domain=35:50,
    samples=100,
    forget plot
]
{
    -(40 + 129.1 * ln(x/35) / ln(10))
};

% Ab 50 MHz
\addplot[
    DARCblue,
    very thick,
    forget plot
]
coordinates {
    (50,-60)
    (1000,-60)
    (1000,-50)
    (3000,-50)
};

\end{axis}

\end{tikzpicture}